\documentclass[11pt]{article}
\usepackage{amsmath,amssymb,amsthm}
\usepackage[margin=1in]{geometry}
\usepackage{booktabs}
\usepackage{array}
\usepackage{tipa}

\newtheorem{defn}{Definition}
\newtheorem{thm}{Theorem}
\newtheorem{prop}{Proposition}
\newcommand{\Mus}{Musta\textcommabelow{t}\u{a}}
\newcommand{\Cont}{\operatorname{Cont}}

\title{Dead-Direction Signatures and Arc Spaces:\\ a Correspondence and its Limits}
\author{}
\date{}

\begin{document}
\maketitle
\thispagestyle{empty}

\section*{Setup}

$K : \mathbb{R}^d \to \mathbb{R}_{\geq 0}$ is the population loss (KL divergence), with
$K(w^*) = 0$, and $W_0 = \{K = 0\}$ is the set of true parameters. $\lambda$ denotes the
RLCT (learning coefficient). For an arc $\gamma$, $\operatorname{ord}_\gamma(K)$ is the
order of vanishing of $t \mapsto K(\gamma(t))$, and
$\Cont^{\geq p}(K) = \{\gamma : \operatorname{ord}_\gamma(K) \geq p\}$, viewed in the jet
space $J_{p-1}(\mathbb{R}^d)$ of dimension $dp$.

\section*{The correspondence}

\begin{defn}[Shirodkar, arXiv:2606.05957]
A unit direction $u$ is a \emph{dead direction of KL order $k$} at $w^*$ if
$K(w^* + tu) = c\,t^{2k} + O(t^{2k+1})$ with $c > 0$.
\end{defn}

\begin{prop}
$u$ is a dead direction of KL order $k$ if and only if the linear arc
$\gamma(t) = w^* + tu$ lies in $\Cont^{\geq 2k}(K) \setminus \Cont^{\geq 2k+1}(K)$.
\end{prop}
\begin{proof}
Both say $\operatorname{ord}_\gamma(K) = 2k$.
\end{proof}

So dead directions are \emph{linear} arcs sorted by contact order. Directions inside
$W_0$ ($K(w^*+tu) \equiv 0$) are not dead directions; they are arcs of infinite order.

\section*{The two formulas for $\lambda$}

\begin{thm}[DDS selection rule; Shirodkar, Thm.~3]
Under that theorem's hypotheses, a dead direction of KL order $k$ contributes
$1/(2k)$ to $\lambda$, and contributions of independent directions add.
\end{thm}

\begin{thm}[\Mus; Ein--Lazarsfeld--\Mus, arXiv:math/0102201, math/0303268]
For an ideal $\mathfrak a$ (or a function), the complex log canonical threshold is
\[
\operatorname{lct}(\mathfrak a) = \min_{p \geq 1}
\frac{\operatorname{codim}_{J_{p-1}} \Cont^{\geq p}(\mathfrak a)}{p}.
\]
\end{thm}

Two conversions connect this to $\lambda$:
\begin{itemize}
\item If $K \asymp f_1^2 + \dots + f_r^2$ near $w^*$, then $\lambda = \operatorname{RLCT}_{\mathbb R}(\mathfrak a)/2$
  for the fiber ideal $\mathfrak a = (f_1,\dots,f_r)$ (Lin, arXiv:1003.5338, Thm.~1.2).
\item $\operatorname{RLCT}_{\mathbb R}(\mathfrak a) = \operatorname{lct}(\mathfrak a)$ when a
  minimizing divisor has enough real points. It can fail: for
  $K = (x^2+y^2+z^2)^2$ the ideal $(x^2+y^2+z^2)$ has $\operatorname{lct} = 1$ but
  $\lambda = 3/4$. Counting \emph{real} arcs avoids this.
\end{itemize}

\section*{How they relate}

DDS is the restriction of the contact formula to linear arcs. The two agree when that
restriction loses nothing, as in the smooth-fibre setting of Thm.~3. The arc formula
also supplies what Thm.~3 takes as given:
\begin{itemize}
\item \textbf{Branch selection}: the minimum over contact-locus components;
\item \textbf{Multiplicity}: the number of minimizing components;
\item \textbf{Non-linear arcs}: a purely linear reading cannot tell a direction tangent
  to a curved $W_0$ from a dead one. For $K = (b - a^2)^2$, $K$ has order 4 along
  $u = (1,0)$, so a linear count gives $1/2 + 1/4$. But the arc $(t, t^2)$ stays in $W_0$,
  so $u$ is tangent, and $\lambda = 1/2$. (Shirodkar's arXiv:2607.00603 classifies such
  directions separately.)
\end{itemize}

\textbf{Where DDS breaks: crossing points.} In reduced-rank regression at the minimum-norm
solution, $W_0$ is a union of components meeting at that point. There
$2\lambda - \operatorname{rank} \nabla^2 K = (H - r_0)(N - r_0)$, and this gap is
\emph{not} higher-order contact. The contact-locus minimum is reached already at $p = 1$,
so $\lambda$ is half the codimension of the largest component of $W_0$. The gap is the
excess of the tangent space over that component (see \texttt{arc/README.md}).
Higher-order (quartic) contact, where the DDS count \emph{is} right, appears instead in
symmetric factorization $WW^\top$.

\section*{Checks}

\begin{enumerate}
\item $K = a^2 b^2$, $u = (1,0)$: $K(tu) \equiv 0$, an arc inside $W_0$ (infinite
  order), not a dead direction.
\item $u = (1,1)/\sqrt 2$: $K(tu) = t^4/4$, KL order 2, contribution $1/4$.
\item Contact formula for $K = a^2b^2$ at $p = 2$:
  $\Cont^{\geq 2} = \{a_0 b_0 = 0\} \subset J_1(\mathbb{R}^2) \cong \mathbb{R}^4$ has
  codimension 1, so the ratio is $1/2$. That equals $\lambda = 1/2$, and also
  $\operatorname{lct}(a^2b^2) = 1/2$. Equivalently, with fiber ideal $(ab)$:
  $\operatorname{lct}(ab) = 1$ and $\lambda = 1/2$. (An earlier version of this note
  computed $2 - 3$ instead of $2 - 3/2$ here and then explained a non-existent
  discrepancy.)
\item Branch selection: on $\{a = 0\}$ and on $\{b = 0\}$ the codimension is 1 and
  $\lambda = 1/2$; at the origin both components are present and the minimum is still
  $1/2$, with multiplicity 2 when the branches have equal codimension.
\end{enumerate}

\end{document}
